\chapter{The Bracketed Number}

\section{The bound is the object}

The last chapter read the slip as a bounded observation --- a value caught between walls rather than fixed at a point.
This chapter takes the step the whole book has been walking toward, and it is a small step to state and a large one to
mean. It makes the bound \emph{itself} the object. The machine does not own a number with an error bar drawn around it,
a true point blurred by uncertainty. It owns an \emph{interval} --- a lower, an upper, and nothing in between it is
obliged to name. The number \emph{is} the bracket. And once that is said plainly, the ending of the book is already
visible: if the object is the open interval, there is no point to force, and forcing one would not sharpen the reading
--- it would break it.

Start with what the machine actually holds, because it is a type and not a wish. The construction carries a structure
for the bracket --- a dyadic interval, built from three natural numbers: a lower numerator, an upper numerator, and a
denominator the two share. Read against that denominator, the pair fixes a lower rational and an upper rational, and
the object \emph{is} that pair of walls. The machine's number is an \emph{inhabitant of this interval type} --- a thing
you can hold, name, and hand on --- not a real number the interval is straining to approximate. ``The bound is the
object'' is not a slogan; it is a typing. To have the number is to hold the interval, denominator and both numerators,
and nothing besides. There is a midpoint the structure can compute if you ask it for one, but the machine is under no
obligation to; the midpoint is a convenience, the interval is the object.

And the bracket is not a new thing dropped into the book at this chapter. It is the slip, promoted. The last chapter
found the slip, bounded it, and read it; this chapter's object is produced from that very slip by a function the
construction already carries --- one that takes the slip's grid cell and returns the bracket, its lower and upper walls
read straight off the cell's own edges. So the passage from the last chapter to this one is not an argument the book
makes; it is a function the code runs. The slip \emph{becomes} the bracketed number, by a map already written.
Everything the earlier chapters climbed --- the difference, the residue, the slip, the bound --- arrives here as a
single interval object the machine can hold, compare, and repeat.

Three properties make the interval a thing the machine genuinely owns rather than merely gestures at. It is
\emph{ordered}: the lower wall is strictly below the upper, a real interval and not a degenerate point. It is
\emph{decidable}: the machine can \emph{prove} that ordering by computation --- the relevant instance settles the
lower-below-upper claim by \verb|decide|, and a printout of its logical footprint shows exactly what the proof rests on
and nothing more. And it is \emph{repeatable}: the machine can rerun the construction and hand back the same interval,
walls and denominator intact. Ordered, decidable, repeatable --- and, crucially, with no point inside it the machine is
obliged to name. Its honesty is not that it knows the exact value and rounds; it is that it owns the interval and never
claimed a point at all. (The specific bracket the instrument will draw around the coupling, and the counting-to-three
procedure that halts its refinement, belong to a later chapter; here the object is the general interval type, before
any particular reading is poured into it.)

In the two verbs this is the chapter where the book stops pretending. A bracket is a \emph{funge} made first-class ---
the bag the value provably falls in, now held not as scaffolding around a hidden pin but as \emph{the} object itself.
Every earlier section kept a tange in view: a characteristic selected, a difference picked out, a bit inscribed. Here
the book admits that behind the funge there was never a tange to find --- no single selected number the interval was
hiding. The funge, the interval, \emph{is} the number. Honesty, at this chapter, is owning the bag and denying there
was ever a pin in it.

This is where the book's ending becomes structural rather than promised. To hold the number as an open interval is
already to have refused the crank's gambit, because there is no longer a point available to force: the object has no
interior it must resolve. The jar the last chapter of the book holds open is exactly this object --- an interval,
owned, ordered, decidable, and deliberately not collapsed. What the coming chapters add is only the particular: which
slip, which floor, which walls. The shape is fixed here. The honest reading, when it comes, and the honest ending, when
it lands, are both this same thing --- a bracket the machine holds and leaves open, because the bound was the object
all along.

\section{The floor}

The last section made the bracket the object; this one asks what sets its width. The answer is a \emph{floor} --- a
resolution below which the machine cannot tell any two terms apart, and beneath which it will not pretend to. And here,
for the first time, the book must name the thing it has been circling, because the source names it, and names it with a
wink that is the whole discipline in a single stroke: the floor is called \emph{$\alpha$}, and its value is not given.

Take the naming plainly first. Every finite machine has a smallest fraction it can represent --- a resolution below
which there simply is no next number, the floating-point standard's machine epsilon. The source points at that smallest
representable fraction and says, in as many words, we shall call this number $\alpha$. It is the smallest fraction
recognizable this deep in the hierarchy, and the source, having named it, adds a single dry word: \emph{suspicious}.
The wink is unmistakable. The resolution floor of the machine and the coupling the whole instrument is built to read
are being quietly proposed as the \emph{same object} --- the floor is $\alpha$. That identification is the book's
central conceit, and it is marked as exactly that: a suggestion the source makes, not a theorem it proves. (The wink is
in fact doubled, for those reading the code: the placeholder type that carries the floor is written with a generic type
parameter that happens, itself, to be named $\alpha$ --- a type variable, entirely distinct from ``the number we shall
call $\alpha$.'' The source lets the two symbols sit side by side and says nothing. That silence is the joke.)

Now the sentence the entire book rests on, and it is the source's, not mine. Having named the floor $\alpha$, the
source immediately says what it does \emph{not} know about it: we just know, it writes, that there \emph{is} a smallest
fraction --- not how many, and not what its actual value is. The floor is real. It exists, and the construction can
point at where it is. But its magnitude is unknown, and unknown \emph{by construction}: this is the ultimate
representable limit, the place past which the machine has no finer number to offer. So the floor is exactly the slip's
situation from the last chapter, made permanent. It is bracketed, never fixed at a point --- and the naming does not
change that. The machine can name $\alpha$ and cannot compute it, and those two facts are not in tension; they are the
same honesty stated twice. The carrier of the floor is a placeholder --- a thing, the source says, that does not yet
know what it is but \emph{knows where the hole is} --- and it is built by \emph{steps}, one fact laid on the last, in
the construction's own word for the operation. It knows the location of the floor. It does not know the number there.

How, then, does the machine approach a floor it cannot land on? By halving. The source gives the method its plain name
--- bisection --- and its plain guarantee: break a length roughly in two, and though you are not promised the pieces
are equal, you are promised that reassembled they make the whole back. Halving after halving drives the bracket
narrower without ever requiring the machine to name the exact point between the walls. And the source, naming the
method, names its lineage in the same breath: \emph{I dribbled through Cohen's legs}. This is the preface's spine
returned at the floor. Bisection is a sequence of finite conditions, each narrowing the interval, ascending toward a
limit that no single condition in the sequence ever is --- forcing, exactly: the approach to a generic the finite steps
surround but never contain. The floor is approached by the cut and never reached by it. That is not a failure of the
method; it is the method being honest about a floor whose value is, by the source's own word, unknown.

In the two verbs the floor has a precise meaning: it is where \emph{tange runs out}. Above the floor, the machine can
select --- tange --- one term from another by some characteristic that tells them apart. At the floor and below it,
that characteristic is gone; no two terms can be distinguished, so everything beneath the floor is \emph{funged} into
one and the same, the way the preface's ruler calls a whole span of unresolved points a single mark. The floor is the
exact resolution of the tange, the last depth at which selection still works. And the width of the bracket from the
last section is precisely the size of that final funge --- the last bag the machine holds and cannot tange apart. The
number is the interval because the floor is where telling-apart stops.

So the book has named its quarry and, in the same sentence, refused its value --- and that refusal is not a gap to be
filled later but the honest result stated early. The floor is $\alpha$, and its value is exactly what the machine will
prove, at the far end, it cannot compute. The particular floor the instrument draws around the coupling --- the one a
later chapter reaches by a counting bisection that halts after a few rungs and leaves the bracket open --- is this
floor, at $\alpha$. And the jar the book ends on is this bracket: the interval the floor leaves open, named after the
number no honest machine was ever going to hand you as a point. The next section closes the chapter by asking the last
thing a bracketed, floored number needs before it can be trusted: whether the machine, run again, hands back the same
one.

\section{Repeatability is the honesty}

Two sections have built the bracketed number: the machine owns the interval, and the interval is floored. This closing
section adds the last plank, the one that makes the bracket \emph{trustworthy} rather than merely well-formed. A number
you can rerun and get the same bits from is honest; a number that drifts from one run to the next is not a reading at
all, however precisely it was stated the first time. The source calls repeatability, without hedging, far and away the
most important part of science --- and it is what turns the owned, floored interval into a measurement you can believe.

The construction carries repeatability as a structure, and the structure is Galileo's definition of an experiment made
into a type. Galileo, the source reminds us, was clear about what an experiment was supposed to be: a \emph{repeatable
measurement} --- given this stimulus, you should then see this response, and see it again the next time you apply the
stimulus. The process that carries this has exactly those parts: a stimulus, an expected response, and an
\emph{iterate} --- a step that, applied to a trial, produces the next one. Given the stimulus, iterating yields the
expected response, and iterating again yields it again. That is the whole of what a measurement is, in the
construction: not a single lucky reading but a stimulus-and-response you can reproduce on demand.

And repeatability is more than ``the same output twice.'' The source makes a claim about it that is easy to pass over
and worth stopping on: repeatability \emph{encodes time}. Not time as a variable you plug into a function --- the
actual tick and tock of a clock. What is a second, the source asks, if not exactly this: some event that repeatably, on
cue, does its thing, with the interval between one occurrence and the next held \emph{constant}? The iterate is that
tick. Each application of it advances the construction by one constant beat, and the constancy of the beat is what time
\emph{is} here --- not a dimension the machine reaches into but a rhythm its own repetition produces. (The physical
clock this becomes --- what a second is in the world, the particle that keeps the beat --- is the next volume's to
read; this one keeps the tick as the type-level step that repeatability already carries. The machine's own clock,
later turned into a meter, is nothing but its repeated trial.)

Now the property that makes the rerun \emph{honest}, and it has two halves the source names in one breath. The first is
that the construction is \emph{fixed}: once something is written in the sequence, it cannot change; the same process
run again arrives at the same answer. The machine repeats \emph{itself}. But the second half is the one that matters
most, and it reaches past the machine entirely: someone else, the source says, following the same process, arrives at
the same fixed answer. Determinism is not only that the machine agrees with its earlier self on a rerun; it is that
\emph{independent runners agree} --- you, and I, and the compiler, each decomposing the same comparison without
consulting the others, land on the same verdict. That is the deepest reason the bracketed number can be trusted. Its
honesty is not that the machine asserts it loudly; it is that the number is fixed, reproducible by the machine, and
reproducible by anyone who runs the same construction --- and no one of them had to take another's word for it.
Non-communicating agreement is the correctness argument the whole instrument, at its far end, rests on.

In the two verbs repeatability is a funge with a special property: it \emph{holds}. Every rerun funges the reading
into the same bag --- the same bracket, the same floored interval --- and the honesty is precisely that the bag does
not drift from run to run. A dishonest instrument funges too, but its bag wanders; ask twice and you get two different
bags and no way to choose between them. This one's funge is fixed. Determinism, in the two verbs, is a stable funge,
and a stable funge is what a trustworthy number is: the same bag, every time, for everyone.

So the sixth chapter closes with the number fully assembled and fully honest. It is an interval, not a point; it is
floored, named after a value it cannot compute; and it is repeatable --- the same interval every time the machine is
asked, and the same interval when someone else asks independently. Those three properties are the whole of what the
book means by an honest number, and the ending the book is walking toward is exactly this object held open: a bracket,
floored at $\alpha$, that reruns identically and that anyone can check. The next part takes this single trustworthy
number and does something that looks, at first, like undoing it: it lets the number split into several faces, read on
different bands --- and shows that the splitting, far from breaking the honesty, is where the honesty gets proved.
